% !TeX root = ../main.tex
\section{Justificación de la arquitectura distribuida}
\label{sec:clab-distribucion}

CLAB parte de una aplicación Angular con backend Spring Boot y persistencia
PostgreSQL. El documento previo identifica reservas coordinadas informalmente,
conflictos de horario, inventario desactualizado e incidencias sin seguimiento.
Estos problemas justifican informatizar el proceso, pero no demuestran por sí
solos que un monolito sea incapaz de resolverlo. No se dispone todavía del acta
firmada o de la grabación consentida que permita verificar el levantamiento
exigido en E1; se usa el documento previo como antecedente, no como evidencia
primaria validada.

La distribución se propone para permitir despliegues por dominio y aislar el
procesamiento de reportes y avisos de la confirmación de reservas. Por ejemplo,
una actualización de estadísticas no debe exigir desplegar de nuevo la agenda;
una caída del servicio de informes no debe deshacer una reserva confirmada.
Un monolito modular es una alternativa viable para la demanda descrita, pero no
ofrece esa independencia de procesos y despliegue. Además, el plan académico
requiere migrar el sistema completo a una base distribuida durante el periodo.
No se usa como justificación una supuesta operación multisede, una saturación
medida o una necesidad de nube pública: ninguna ha sido comprobada.

Los requisitos E1 actuales concretan las metas del antecedente. RNF-01 propone
que al menos el 95\,\% de las consultas válidas de disponibilidad entregue un
resultado correcto en menos de 2 segundos, durante 10 minutos de carga estable
con 50 usuarios concurrentes. Errores y tiempos agotados cuentan como
incumplimientos. RNF-02 propone que al menos el 99\,\% de las comprobaciones de
una sonda por minuto sea correcto dentro de 5 segundos, durante dos semanas
de horario académico acordado. Son umbrales propuestos, no resultados medidos;
la configuración concreta del entorno y su ejecución deben registrarse.

Separar servicios podría permitir escalar consultas e informes sin replicar
toda la aplicación, pero agrega llamadas de red y no prueba esas metas.
La disponibilidad se evalúa por operación: una caída de Avisos no debe impedir
persistir reservas válidas, mientras que perder Identidad impide autorizar
solicitudes protegidas. Un único servidor o una única instancia PostgreSQL
siguen siendo puntos comunes de fallo.

\subsection{Trazabilidad hacia los requisitos E1 actuales}

Se usan los identificadores de la propuesta funcional de \texttt{main.tex}
y los RNF de \texttt{Carlos.tex}, consultados en la revisión
\texttt{a2ecc7d}. No se confunden con los identificadores homónimos del informe
previo, ni se presentan como requisitos ya aprobados o implementados:
\begin{description}
  \item[RF-P01 a RF-P03:] la consulta, solicitud y cancelación se apoyan en
  Reservas como autoridad de ocupación. Laboratorios entrega metadatos y
  estado; Identidad verifica permisos antes de ejecutar cambios.
  \item[RF-P04:] el catálogo pertenece a Laboratorios. Un estado fuera de
  servicio detectado al validar impide una nueva confirmación. Aún debe
  acordarse el instante efectivo del cierre si la desactivación concurre con
  una reserva: consultar el catálogo no hace atómica una operación entre dos
  almacenes. Este punto se registra como pendiente, no como regla ya resuelta.
  \item[RNF-01 y RNF-02:] orientan pruebas de latencia y disponibilidad sobre
  consultas de agenda. La separación permite evaluar el escalado selectivo,
  pero no elimina la dependencia de Identidad o de infraestructura compartida.
  \item[RNF-03:] confirmar, cancelar o bloquear requiere autorización
  verificable. Una revocación o pérdida de Identidad no se sustituye por
  asumir que todo token todavía vigente conserva permisos.
  \item[RNF-04 y RNF-05:] la SPA debe soportar los flujos y tamaños de pantalla
  especificados; los servicios deben poder desplegarse sin nube pública.
  Distribuir no obliga a contratar infraestructura externa.
  \item[RNF-06:] diez solicitudes válidas de la prueba, para franjas libres
  diferentes, deben confirmarse y persistirse con Información y Avisos
  detenido, manteniendo las dependencias de agenda operativas. Los avisos
  quedan fuera de la transacción crítica de confirmación.
  \item[RNF-07:] cada cambio exitoso de agenda del conjunto de prueba requiere
  auditoría en su servicio de origen, incluso si la vista central está caída.
  La copia diferida no sustituye ese registro local.
\end{description}

\subsection{Resumen de ADR-001}

\textbf{Estado: propuesta pendiente de aceptación del equipo.} Se compararon
un monolito modular, doce servicios ---uno por módulo--- y cuatro servicios por
dominio. Se propone la tercera opción para limitar la complejidad operativa:
\begin{itemize}
  \item \textbf{Identidad y Académico:} usuarios, roles, permisos, datos
  académicos y estudiantes.
  \item \textbf{Laboratorios y Recursos:} catálogo, equipos e incidencias.
  \item \textbf{Reservas y Agenda:} reservas, bloqueos y asistencia.
  \item \textbf{Información y Avisos:} reportes de uso, notificaciones y vista
  central de auditoría; la auditoría original permanece en cada dominio.
\end{itemize}

Así se cubren los doce módulos sin exigir doce despliegues. Cada servicio es
dueño de su almacenamiento; no accede a tablas de otro servicio. Reservas
conserva juntas las reservas y los bloqueos para decidir la ocupación en una
transacción local. Se propone que todas las mutaciones de agenda adquieran el
mismo bloqueo de fila por laboratorio antes de comprobar solapamientos y
persistir. La documentación de PostgreSQL distingue restricciones entre filas
de simples comprobaciones locales; una consulta previa aislada no sustituye
la protección concurrente\cite{clabPostgresRestricciones}.

La comunicación REST/JSON ya existe en el esqueleto de E1. El diseño futuro
propone consultas sincrónicas para autorización y validación del catálogo,
y lectura incremental de cambios durables en segundo plano para informes y
avisos. Los contratos y parámetros deben cerrarse en ADR-002 con los resultados
reales de PE-U1, que no estuvieron disponibles en esta revisión. No se atribuye
a REST una superioridad experimental todavía no demostrada.

La migración prevista es incremental: E1 define límites y demuestra dos procesos
comunicándose; E2 incorpora APIs y contenedores; E3 completa la separación de
datos, los módulos restantes y las proyecciones; E4 valida la migración completa,
la recuperación y la calidad. Cada corte mantiene una única fuente de escritura
por entidad y verifica los procedimientos SQL del legado antes de sustituirlos.
El registro completo de opciones, consecuencias y pruebas pendientes está en
\texttt{docs/adr/ADR-001-estilo-arquitectonico.md}.

\paragraph{Estado de implementación y autoría.}
En la revisión del código se identificaron dos servicios mínimos sin persistencia
de negocio. Los cuatro servicios y sus garantías son diseño propuesto, no
integración realizada. Esta sección y ADR-001 se prepararon con asistencia de
IA (OpenAI, GPT-6.1 Sol); no acreditan validación institucional ni aprobación
del equipo.
